\documentclass[10pt,a4paper]{amsart}
\usepackage{amsmath,amssymb,amsthm,graphicx}
\usepackage{verbatim}
\newenvironment{reference}{\comment}{\endcomment}
%\newenvironment{reference}{}{}
\newenvironment{slogan}{\comment}{\endcomment}
\newenvironment{history}{\comment}{\endcomment}

% For commutative diagrams we use Xy-pic
\usepackage[all]{xy}

% We use 2cell for 2-commutative diagrams.
\xyoption{2cell}
% Theorem environments.
%
\theoremstyle{plain}
\newtheorem{theorem}[subsection]{Theorem}
\newtheorem{proposition}[subsection]{Proposition}
\newtheorem{lemma}[subsection]{Lemma}

\theoremstyle{definition}
\newtheorem{definition}[subsection]{Definition}
\newtheorem{example}[subsection]{Example}
\newtheorem{exercise}[subsection]{Exercise}
\newtheorem{situation}[subsection]{Situation}

\theoremstyle{remark}
\newtheorem{remark}[subsection]{Remark}
\newtheorem{remarks}[subsection]{Remarks}

\numberwithin{equation}{subsection}

% Macros
%
\def\lim{\mathop{\mathrm{lim}}\nolimits}
\def\colim{\mathop{\mathrm{colim}}\nolimits}
\def\Spec{\mathop{\mathrm{Spec}}}
\def\Hom{\mathop{\mathrm{Hom}}\nolimits}
\def\Ext{\mathop{\mathrm{Ext}}\nolimits}
\def\SheafHom{\mathop{\mathcal{H}\!\mathit{om}}\nolimits}
\def\SheafExt{\mathop{\mathcal{E}\!\mathit{xt}}\nolimits}
\def\Sch{\mathit{Sch}}
\def\Mor{\mathop{\mathrm{Mor}}\nolimits}
\def\Ob{\mathop{\mathrm{Ob}}\nolimits}
\def\Sh{\mathop{\mathit{Sh}}\nolimits}
\def\NL{\mathop{N\!L}\nolimits}
\def\CH{\mathop{\mathrm{CH}}\nolimits}
\def\proetale{{pro\text{-}\acute{e}tale}}
\def\etale{{\acute{e}tale}}
\def\QCoh{\mathit{QCoh}}
\def\Ker{\mathop{\mathrm{Ker}}}
\def\Im{\mathop{\mathrm{Im}}}
\def\Coker{\mathop{\mathrm{Coker}}}
\def\Coim{\mathop{\mathrm{Coim}}}

% Boxtimes
%
\DeclareMathSymbol{\boxtimes}{\mathbin}{AMSa}{"02}

%
% Macros for moduli stacks/spaces
%
\def\QCohstack{\mathcal{QC}\!\mathit{oh}}
\def\Cohstack{\mathcal{C}\!\mathit{oh}}
\def\Spacesstack{\mathcal{S}\!\mathit{paces}}
\def\Quotfunctor{\mathrm{Quot}}
\def\Hilbfunctor{\mathrm{Hilb}}
\def\Curvesstack{\mathcal{C}\!\mathit{urves}}
\def\Polarizedstack{\mathcal{P}\!\mathit{olarized}}
\def\Complexesstack{\mathcal{C}\!\mathit{omplexes}}
% \Pic is the operator that assigns to X its picard group, usage \Pic(X)
% \Picardstack_{X/B} denotes the Picard stack of X over B
% \Picardfunctor_{X/B} denotes the Picard functor of X over B
\def\Pic{\mathop{\mathrm{Pic}}\nolimits}
\def\Picardstack{\mathcal{P}\!\mathit{ic}}
\def\Picardfunctor{\mathrm{Pic}}
\def\Deformationcategory{\mathcal{D}\!\mathit{ef}}
\usepackage{cleveref}
\crefname{originallemma}{original lemma}{original lemmas}
\Crefname{originallemma}{Original lemma}{Original lemmas}
\makeatletter
\newlabel{more-morphisms-lemma-go-down-with-annihilators}{{082V}{source}{}{sourcepointer-1}{}}
\newlabel{more-morphisms-lemma-go-down-with-annihilators@cref}{{[originallemma][0][]082V}{[1][0][]source}{}{}{}}
\makeatother
\crefname{originaldefinition}{original definition}{original definitions}
\Crefname{originaldefinition}{Original definition}{Original definitions}
\makeatletter
\newlabel{morphisms-definition-scheme-theoretically-dense}{{01RB}{source}{}{sourcepointer-2}{}}
\newlabel{morphisms-definition-scheme-theoretically-dense@cref}{{[originaldefinition][0][]01RB}{[1][0][]source}{}{}{}}
\makeatother
\crefname{originallemma}{original lemma}{original lemmas}
\Crefname{originallemma}{Original lemma}{Original lemmas}
\makeatletter
\newlabel{morphisms-lemma-characterize-scheme-theoretically-dense}{{01RE}{source}{}{sourcepointer-3}{}}
\newlabel{morphisms-lemma-characterize-scheme-theoretically-dense@cref}{{[originallemma][0][]01RE}{[1][0][]source}{}{}{}}
\makeatother
\title{The preimage of a scheme theoretically dense open under a flat, locally finitely presented, locally quasi-finite morphism is scheme theoretically dense}
\author{The Stacks Project Authors}
\date{stacks-project:spaces-morphisms.tex:local-snapshot-20261012; GNU FDL 1.2 or later}
\begin{document}
\maketitle
\setcounter{subsection}{1}
\noindent\textit{Editorial note.} All delivered mathematical text is the author's own, copied from the pinned Stacks Project chapter file spaces-morphisms.tex; no statement or proof step was rewritten, reordered or removed.\par\medskip
\noindent\textit{Editorial note.} The statement is about morphisms of schemes, and its proof is likewise scheme-theoretic; it is a chapter tool used later for algebraic spaces and is counted as such.\par\medskip
\noindent\textit{Editorial note.} The three references of the counted spans are recorded as non-load-bearing navigation pointers; the argument that consumes them stays inside the counted proof span.\par\medskip
\noindent\textit{Editorial note.} Preamble: the three native excerpt spans of sources/stacks-spaces-morphisms/preamble.tex (lines 8-13, 14-18 and 154-217). The chapter file itself carries no preamble text.\par\medskip
\noindent\textit{Allowed standard prerequisites (editorial).}\begin{itemize}
\item Characterization of scheme theoretic density by functions vanishing on the open. Morphisms, Lemma morphisms-lemma-characterize-scheme-theoretically-dense, used at proof lines 3432-3436. Reduces density to the statement that such a function g must vanish.
\item The go-down-with-annihilators lemma for flat locally finitely presented locally quasi-finite morphisms. More on Morphisms, Lemma more-morphisms-lemma-go-down-with-annihilators, used at proof lines 3437-3450. Produces U, the subsheaf F, the integer r and the injection carrying g.
\item Scheme theoretic density of W in S. used at proof lines 3452-3454. Forces the sections g\_i to vanish.
\end{itemize}\medskip
\ifdefined\hypertarget\hypertarget{sourcepointer-1}{}\fi
\noindent\textit{Source reference.} Original lemma 082V: focused-seeds/AGP/sources/stacks-snapshot-20261011/tags.txt, source lines 10265--10265. Navigation pointer into the Stacks Project; the standard fact it supplies is applied inside the counted proof span and the argument resting on it stays there.\par\medskip
\ifdefined\hypertarget\hypertarget{sourcepointer-2}{}\fi
\noindent\textit{Source reference.} Original definition 01RB: focused-seeds/AGP/sources/stacks-snapshot-20261011/tags.txt, source lines 2296--2296. Navigation pointer into the Stacks Project; the standard fact it supplies is applied inside the counted proof span and the argument resting on it stays there.\par\medskip
\ifdefined\hypertarget\hypertarget{sourcepointer-3}{}\fi
\noindent\textit{Source reference.} Original lemma 01RE: focused-seeds/AGP/sources/stacks-snapshot-20261011/tags.txt, source lines 2299--2299. Navigation pointer into the Stacks Project; the standard fact it supplies is applied inside the counted proof span and the argument resting on it stays there.\par\medskip
\begin{lemma}
\label{lemma-scheme-theoretically-dense-representable}
Let $S$ be a scheme. Let $W \subset S$ be a scheme theoretically
dense open subscheme
(Morphisms, Definition \ref{morphisms-definition-scheme-theoretically-dense}).
Let $f : X \to S$ be a morphism of schemes which is flat, locally of
finite presentation, and locally quasi-finite.
Then $f^{-1}(W)$ is scheme theoretically dense in $X$.
\end{lemma}

\begin{proof}
We will use the characterization of Morphisms, Lemma
\ref{morphisms-lemma-characterize-scheme-theoretically-dense}.
Assume $V \subset X$ is an open and $g \in \Gamma(V, \mathcal{O}_V)$
is a function which restricts to zero on $f^{-1}(W) \cap V$.
We have to show that $g = 0$. Assume $g \not = 0$ to get a
contradiction. By
More on Morphisms, Lemma \ref{more-morphisms-lemma-go-down-with-annihilators}
we may shrink $V$, find an open $U \subset S$ fitting into a
commutative diagram
$$
\xymatrix{
V \ar[r] \ar[d]_\pi & X \ar[d]^f \\
U \ar[r] & S,
}
$$
a quasi-coherent subsheaf $\mathcal{F} \subset \mathcal{O}_U$, an integer
$r > 0$, and an injective $\mathcal{O}_U$-module map
$\mathcal{F}^{\oplus r} \to \pi_*\mathcal{O}_V$
whose image contains $g|_V$. Say
$(g_1, \ldots, g_r) \in \Gamma(U, \mathcal{F}^{\oplus r})$ maps to $g$.
Then we see that $g_i|_{W \cap U} = 0$ because $g|_{f^{-1}W \cap V} = 0$.
Hence $g_i = 0$ because $\mathcal{F} \subset \mathcal{O}_U$ and
$W$ is scheme theoretically dense in $S$.
This implies $g = 0$ which is the desired contradiction.
\end{proof}

\end{document}
